\documentclass[11pt]{article}

% ----------------------------------------------------------------------
%  The Synergy Reality Test — client-facing methodology one-pager
%  Companion to value-creation-bridge.tex. Draft concept by the ENT 105
%  team for Creo Advisors. Compile: pdflatex synergy-reality-test.tex
% ----------------------------------------------------------------------

\usepackage[utf8]{inputenc}
\usepackage[T1]{fontenc}
\usepackage{textcomp}
\usepackage[margin=0.9in]{geometry}
\usepackage{booktabs}
\usepackage{array}
\usepackage{enumitem}
\usepackage{titlesec}
\usepackage{xcolor}
\usepackage{tcolorbox}
\usepackage[hidelinks]{hyperref}

\definecolor{creonavy}{RGB}{17,38,73}
\definecolor{creoaccent}{RGB}{27,94,140}

\titleformat{\section}{\large\bfseries\color{creonavy}}{}{0em}{}
\titlespacing*{\section}{0pt}{1.5ex plus .3ex}{0.7ex}
\setlist[itemize]{leftmargin=1.3em,topsep=2pt,itemsep=2.5pt}
\setlength{\parskip}{4pt plus 1pt minus 1pt}
\setlength{\parindent}{0pt}

\newtcolorbox{hook}{colback=creonavy!6,colframe=creonavy,
  boxrule=0.8pt,arc=2pt,left=8pt,right=8pt,top=6pt,bottom=6pt}
\newtcolorbox{sowhat}{colback=creoaccent!8,colframe=creoaccent,
  boxrule=0.6pt,arc=2pt,left=8pt,right=8pt,top=5pt,bottom=5pt}

\newcommand{\durable}{\textcolor{creoaccent}{\textbf{Durable}}}
\newcommand{\fragile}{\textcolor{gray!55!black}{\textbf{Fragile}}}

\begin{document}

\begin{center}
{\LARGE\bfseries\color{creonavy} The Synergy Reality Test}\\[4pt]
{\large\color{creoaccent} How much of your roll-up's synergy is durable --- and how much
will competition take back?}\\[5pt]
{\normalsize Separating real efficiency from fragile rent in buy-and-build.}
\end{center}

\vspace{2pt}\hrule\vspace{8pt}

\begin{hook}
\textbf{Margin going up is not the question.} In a buy-and-build, consolidating the market
\emph{always} lifts pricing power --- so a synergy model that shows margins rising proves
nothing. The total rent in a market is conserved; consolidation only \textbf{moves it from the
pool competition was eroding into the pool you now protect}. A roll-up \textbf{captures} rent ---
it does not create it. The synergy is real only to the extent it is \textbf{durable}. The Synergy
Reality Test tells you how much of it is.
\end{hook}

\section{What it answers}
\begin{enumerate}[leftmargin=1.6em]
  \item \textbf{Durable or fragile?} Of the projected synergy, how much is \durable{} ---
  genuine efficiency you recaptured, which sticks --- versus \fragile{}: margin gained
  simply by having fewer competitors, which a new entrant can undo and which draws antitrust
  scrutiny?
  \item \textbf{How far can you consolidate?} Where is the \textbf{natural floor} --- the
  free-entry level of players below which the extra margin sits above what entry will tolerate,
  and gets competed back away?
\end{enumerate}

\section{How it works}
\textbf{Capture, not creation.} The total economic rent in a market is conserved; consolidation
only moves it between two pools --- the rent competition is dissipating, and the rent you
protect. A roll-up shifts rent from the first into the second; the total is unchanged. That
reframes ``synergy'' as a split, not a sum:
\begin{itemize}
  \item \durable{} \textbf{share --- efficiency recaptured.} Competition wastes resources
  (duplicated overhead, value-destroying price wars). Consolidating eliminates that waste ---
  a genuine, lasting gain that takes nothing from anyone.
  \item \fragile{} \textbf{share --- competition-reduction rent.} The rest is margin you earn
  purely by being one of fewer players. It is a transfer from customers, it invites new
  entrants, and it is the part antitrust cares about. It does not last.
\end{itemize}
\textbf{The natural floor.} Every market has a free-entry number of players, set by how costly it
is to enter and how differentiated the product is. Consolidate \emph{to} that number and the gains
are durable; push \emph{below} it and the extra margin rises above what entry tolerates and gets
competed back. There is a markup entry cannot erode --- the durable ceiling on what consolidation
can hold.

\section{Why you can trust it --- and what it won't tell you}
\textbf{What's rigorous:} the test rests on the same conservation backbone as the Value-Creation
Bridge --- because total rent is conserved, captured rent splits into durable and fragile with no
overlap and no residual, and the natural floor is the market's free-entry level, a formal result
rather than a rule of thumb. \textbf{The split arithmetic and the free-entry-floor formula are
exact and self-checking} --- you can verify them on the page; the mechanics, not the inputs.

\smallskip
\textbf{What it deliberately does \emph{not} claim:} it is \textbf{not} an ``optimal number of
acquisitions'' calculator. On its own the geometry says ``consolidate fully or not at all''; a real
stopping point depends on integration costs and antitrust risk, which we factor in as
\emph{explicit, stated} inputs --- not something the math decides for you. And it gives the durable
\emph{level}, not the \emph{speed} at which fragile rent erodes; how fast entry arrives is a
judgment we flag, never a number we invent.

\section{What we need / what you get}
\begin{center}
\renewcommand{\arraystretch}{1.25}\small
\begin{tabular}{p{6.6cm} p{6.6cm}}
\toprule
\textbf{Inputs (from you, stated as assumptions)} & \textbf{Outputs (from us)} \\
\midrule
An estimate of the durable share (genuine efficiency vs.\ competition reduction); entry cost and
product differentiation; antitrust thresholds. &
The synergy split into durable vs.\ fragile; the free-entry floor and durable markup; a
``price-on-the-durable-number'' verdict; and a flag on which acquisitions buy fragile rent. \\
\bottomrule
\end{tabular}
\end{center}

% ======================================================================
\section{Worked example (illustrative)}
% ======================================================================
A platform in a fragmented market of \textbf{6 competitors} pitches a buy-and-build to
consolidate to \textbf{3}, projecting \textbf{\$60M} of annual EBITDA synergy --- which, in full,
is rent the consolidation \emph{captures}. The Test runs two checks.

\smallskip
\textbf{Check 1 --- of the \$60M, how much is durable?} We estimate the durable share at
\textbf{60\%} --- genuine cost/efficiency recapture rather than pure competition reduction --- and
apply it to the captured \$60M:
\begin{center}
\renewcommand{\arraystretch}{1.25}\small
\begin{tabular}{l r l}
\toprule
\textbf{Projected synergy} & \textbf{\$M} & \textbf{Character} \\
\midrule
Efficiency genuinely recaptured (durable, 60\%) & 36 & \durable{} --- sticks \\
Competition-reduction rent (40\%)               & 24 & \fragile{} --- reverts \\
\midrule
\textbf{Headline synergy} & \textbf{60} & \\
\bottomrule
\end{tabular}
\end{center}

\textbf{Check 2 --- how far can you consolidate before it's fragile?} The market's free-entry
floor works out to \textbf{4 players}, with an entry-proof markup entry cannot erode below:
\begin{center}
\renewcommand{\arraystretch}{1.25}\small
\begin{tabular}{l l}
\toprule
\textbf{Consolidation step} & \textbf{Character} \\
\midrule
6 $\rightarrow$ 4 players & \durable{} --- consolidating down to the natural floor \\
4 $\rightarrow$ 3 players & \fragile{} --- one step below the floor; entry restores $\sim$4 and competes the margin away \\
\bottomrule
\end{tabular}
\end{center}

{\footnotesize The figures are the exact decomposition, with no residual. \emph{Check~1:} with a
durable share $\rho=0.60$ on the captured \$60M, the durable part is $\rho\times 60=36$ and the
fragile part $(1-\rho)\times 60=24$ (\$M), reconciling to \$60M. \emph{Check~2} is structural ---
not a second dollar deduction: the natural floor is the free-entry count
$n^{*}=\sqrt{t/f}=\sqrt{20/1.25}=4$ players (differentiation $t$ and entry cost $f$ taken as \$20M
and \$1.25M), at which the entry-proof markup is $\sqrt{tf}=5$ (\$M) and durable concentration is
$1/n^{*}=\sqrt{f/t}=0.25$. The two checks are complementary lenses on the same fragile rent --- the
$\rho$-split sizes it in dollars, the floor locates it in the deal sequence --- and are never added.}

\begin{sowhat}
\textbf{So what.} The \$60M pitch is really \textbf{\$36M of durable synergy.} The other
\textbf{\$24M is fragile} --- competition-reduction rent that entry and antitrust will erode. Check
2 locates it: the last acquisition (\mbox{$4\!\rightarrow\!3$}) pushes below the natural 4-player
floor, so its margin sits above the level entry tolerates. \textbf{Price the platform on the
durable \$36M, not the headline \$60M; treat 4 players as the natural floor; and recognize the last
deal buys rent that reverts.} If you still pursue it, do so with eyes open --- and budget for the
entry and antitrust response.
\end{sowhat}

\textbf{Assumptions --- stated, because that is the point.} The durable share (60\%) is your
estimate of real efficiency versus competition reduction; the natural floor (4 players) follows
from the market's entry cost and differentiation; reversion is treated as eventual, with timing
flagged as a judgment, not modeled. The split arithmetic and the free-entry-floor formula are
fixed; the inputs are yours --- change any and the test updates transparently. The rigor is exactly
as good as the assumptions, and the assumptions are visible.

\vspace{6pt}\hrule\vspace{5pt}
{\footnotesize\itshape Companion to ``The Value-Creation Bridge.'' Draft methodology concept
prepared by the ENT 105 team for Creo Advisors. Figures are illustrative. The durable/fragile split
and the free-entry floor are grounded in conservation-based accounting; the durable share, entry cost, differentiation, antitrust thresholds, and
reversion timing are client-specific inputs set at the engagement's scoping stage.}

\end{document}
